% Referencias bibliográficas (APA 7) — desde docs/informe-final/referencias.md
\chapter{Referencias}

Las referencias se presentan en formato APA 7.\textsuperscript{a} edición, ordenadas alfabéticamente por apellido del primer autor.

\vspace{1em}
\begingroup
\setlength{\parindent}{-1.5em}
\setlength{\leftskip}{1.5em}
\setlength{\parskip}{0.6em}

Amid, A., Moalagh, M., \& Ravasan, A. Z. (2012). Identification and classification of ERP critical
failure factors in Iranian industries. \emph{Information Systems}, \emph{37}(3), 227--237.
\url{https://doi.org/10.1016/j.is.2011.10.010}

Banco Interamericano de Desarrollo. (2021). \emph{La transformación digital de las pymes latinoamericanas:
Diagnóstico y hoja de ruta}. BID. \url{https://publications.iadb.org/}

Bangor, A., Kortum, P. T., \& Miller, J. T. (2008). An empirical evaluation of the System Usability
Scale. \emph{International Journal of Human-Computer Interaction}, \emph{24}(6), 574--594.
\url{https://doi.org/10.1080/10447310802205776}

Benlian, A., \& Hess, T. (2011). Opportunities and risks of software-as-a-service: Findings from a
survey of IT executives. \emph{Decision Support Systems}, \emph{52}(1), 232--246.
\url{https://doi.org/10.1016/j.dss.2011.07.007}

Bezemer, C. P., \& Zaidman, A. (2010). Multi-tenant SaaS applications: Maintenance dream or
nightmare? En \emph{Proceedings of the Joint ERCIM Workshop on Software Evolution (EVOL) and
International Workshop on Principles of Software Evolution (IWPSE)} (pp.~88--92). ACM.
\url{https://doi.org/10.1145/1862372.1862393}

Braun, V., \& Clarke, V. (2006). Using thematic analysis in psychology. \emph{Qualitative Research in
Psychology}, \emph{3}(2), 77--101. \url{https://doi.org/10.1191/1478088706qp063oa}

Brooke, J. (1996). SUS: A ``quick and dirty'' usability scale. En P. W. Jordan, B. Thomas,
B. A. Weerdmeester \& I. L. McClelland (Eds.), \emph{Usability evaluation in industry} (pp.~189--194).
Taylor \& Francis.

Campbell, D. T., \& Stanley, J. C. (1963). \emph{Experimental and quasi-experimental designs for
research}. Rand McNally.

Cardona, M., Vera, C., \& Tabares, J. (2022). Factores críticos en la implementación de sistemas ERP
en pymes colombianas: Una revisión sistemática de la literatura. \emph{Información Tecnológica},
\emph{33}(1), 187--196. \url{https://doi.org/10.4067/S0718-07642022000100187}

Comisión Económica para América Latina y el Caribe. (2022). \emph{Perspectivas económicas de América
Latina 2022: Hacia una transición verde y digital}. CEPAL.
\url{https://repositorio.cepal.org/}

Creswell, J. W., \& Plano Clark, V. L. (2018). \emph{Designing and conducting mixed methods research}
(3rd ed.). SAGE.

Duan, J., Faker, K., Fesak, A., \& Stuart, T. (2012). Benefits and drawbacks of cloud-based versus
traditional ERP systems. En \emph{Proceedings of Advanced Topics in Information Technology}
(pp.~1--12).

Ferraiolo, D. F., Sandhu, R. S., Gavrila, S., Kuhn, D. R., \& Chandramouli, R. (2001). Proposed NIST
standard for role-based access control. \emph{ACM Transactions on Information and System Security},
\emph{4}(3), 224--274. \url{https://doi.org/10.1145/501978.501980}

Fielding, R. T. (2000). \emph{Architectural styles and the design of network-based software architectures}
{[}Tesis doctoral, University of California, Irvine{]}.
\url{https://www.ics.uci.edu/\textasciitilde} fielding/pubs/dissertation/top.htm

Fowler, M., \& Foemmel, M. (2006). \emph{Continuous integration}. ThoughtWorks.
\url{https://www.thoughtworks.com/continuous-integration}

Gackenheimer, C. (2015). \emph{Introduction to React}. Apress.

Guo, C., Zhang, J., \& Liu, B. (2017). A SaaS multi-tenancy performance evaluation approach based on
isolation model. En \emph{Proceedings of the 2017 IEEE 4th International Conference on Cyber Security
and Cloud Computing (CSCloud)} (pp.~277--282). IEEE.
\url{https://doi.org/10.1109/CSCloud.2017.61}

Klaus, H., Rosemann, M., \& Gable, G. G. (2000). What is ERP? \emph{Information Systems Frontiers},
\emph{2}(2), 141--162. \url{https://doi.org/10.1023/A:1026543906354}

Kumar, K., \& van Hillegersberg, J. (2000). ERP experiences and evolution. \emph{Communications of the
ACM}, \emph{43}(4), 22--26. \url{https://doi.org/10.1145/332051.332063}

Ramírez, S. (2018). \emph{FastAPI: High performance, easy to learn, fast to code, ready for production}.
\url{https://fastapi.tiangolo.com/}

Ramírez, S. (2021). \emph{SQLModel: SQL databases in Python, designed for simplicity, compatibility, and
robustness}. \url{https://sqlmodel.tiangolo.com/}

Reichheld, F. F. (2003). The one number you need to grow. \emph{Harvard Business Review}, \emph{81}(12),
46--54.

Rodríguez-Abitia, G., Bribiesca-Correa, G., García-Rodríguez, F. J., \& Martínez-Alonso, M. (2020).
Digital gap in universities and challenges for digital transformation in Mexico during COVID-19.
\emph{Future Internet}, \emph{12}(12), Artículo 210. \url{https://doi.org/10.3390/fi12120210}

Soh, C., Kien, S. S., \& Tay-Yap, J. (2000). Cultural fits and misfits: Is ERP a universal solution?
\emph{Communications of the ACM}, \emph{43}(4), 47--51. \url{https://doi.org/10.1145/332051.332070}

Yin, R. K. (2018). \emph{Case study research and applications: Design and methods} (6th ed.). SAGE.

Zhang, Z., Li, X., \& Wang, C. (2021). Usability evaluation of ERP systems for SME users:
A multi-criteria approach. \emph{Computers in Human Behavior}, \emph{118}, Artículo 106698.
\url{https://doi.org/10.1016/j.chb.2020.106698}

\par\endgroup

\vspace{1em}
\begin{quote}
\textbf{Nota}: Las versiones exactas de todas las dependencias de software utilizadas en el proyecto se especifican en los archivos \texttt{backend/pyproject.toml} y \texttt{frontend/package.json} del repositorio.
\end{quote}
